\documentclass[10pt]{article}
\usepackage[T1]{fontenc}
\usepackage[utf8]{inputenc}
\usepackage[french]{babel}
\usepackage{amsmath,amssymb,mathrsfs,mathpazo}
\usepackage{geometry}
\geometry{a4paper,margin=2cm}
\usepackage{graphicx,array,tabularx,tabularray,colortbl,multirow,multicol,booktabs,enumitem,pifont,ulem,fancybox,diagbox,cancel}
\usepackage{pstricks,pstricks-add,pst-plot,pst-tree,pst-node,pst-text,pst-func,pst-math,pst-bspline,pst-3dplot,pst-eucl,pst-solides3d}
\let\pscalculate\undefined
\usepackage{pst-bezier}
\usepackage{tikz,pgfplots,tkz-tab}
\usetikzlibrary{arrows,arrows.meta,patterns,calc,positioning,trees,decorations.markings,shapes.geometric,shapes.misc,shapes.arrows,intersections,decorations.pathmorphing,decorations.pathreplacing,backgrounds,shadows,babel}
\pgfplotsset{compat=1.18}
\usepgfplotslibrary{fillbetween}
\usepackage[np]{numprint}
\usepackage{listings,makecell,pas-tableur,twoopt}
\let\barre\undefined
\usepackage[pstricks]{tablvar}
\definecolor{darkgreen}{RGB}{0,100,0}
\providecommand{\red}{\color{red}}
\providecommand{\blue}{\color{blue}}
\providecommand{\green}{\color{green}}
\providecommand{\black}{\color{black}}
\providecommand{\white}{\color{white}}
\providecommand{\gray}{\color{gray}}
\providecommand{\cyan}{\color{cyan}}
\providecommand{\magenta}{\color{magenta}}
\providecommand{\yellow}{\color{yellow}}
\providecommand{\violet}{\color{violet}}
\providecommand{\orange}{\color{orange}}
\providecommand{\euro}{\texteuro}
\providecommand{\EUR}{\texteuro}
\providecommand{\ieme}{\textsuperscript{e}}
\providecommand{\degres}{\ensuremath{^\circ}}
\providecommand{\degree}{\ensuremath{^\circ}}
\providecommand{\vv}[1]{\overrightarrow{#1}}
\providecommand{\Coord}[1]{#1}
\providecommand{\decosix}{}
\providecommand{\starredbullet}{\ensuremath{\bullet}}
\pagestyle{empty}
\setlength{\parindent}{0pt}
\usepackage{adjustbox}\DeclareUnicodeCharacter{2236}{:}\DeclareUnicodeCharacter{2190}{\ensuremath{\leftarrow}}
\providecommand{\R}{\mathbb{R}}
\providecommand{\N}{\mathbb{N}}
\providecommand{\Z}{\mathbb{Z}}
\providecommand{\Q}{\mathbb{Q}}
\providecommand{\D}{\mathbb{D}}
\providecommand{\C}{\mathbb{C}}
\providecommand{\e}{{\mathrm{e}}}
\providecommand{\vect}[1]{\overrightarrow{#1}}
\providecommand{\vectt}[1]{\overrightarrow{\mathrm{#1}}}
\providecommand{\vvt}[1]{\overrightarrow{\mathrm{#1}}}
\providecommand{\np}[2][]{{#2~#1}}
\providecommand{\vv}[1]{{\overrightarrow{#1}}}
\renewcommand{\d}{{\mathrm{d}}}
\providecommand{\degres}{{^\circ}}
\providecommand{\degre}{{^\circ}}
\providecommand{\blue}{\color{blue}}
\providecommand{\red}{\color{red}}
\providecommand{\enskip}{\quad}
\providecommand{\llbracket}{\lbrack\!\lbrack}
\providecommand{\rrbracket}{\rbrack\!\rbrack}
\providecommand{\euro}{{\text{€}}}
\providecommand{\no}{{\text{n°}}}
\providecommand{\interval}[2]{{[#1;#2]}}
\providecommand{\intervFO}[2]{{[#1;#2[}}
\providecommand{\dotfill}{\ldots}

\providecommand{\ds}{\displaystyle}
\providecommand{\textding}[1]{#1}
\providecommand{\Oij}{\(\mathrm{O};\vec{\imath},\vec{\jmath}\)}
\providecommand{\Oijk}{\(\mathrm{O};\vec{\imath},\vec{\jmath},\vec{k}\)}
\providecommand{\Ouv}{\(\mathrm{O};\vec{u},\vec{v}\)}
\providecommand{\intervFF}[2]{\ensuremath{[#1;#2]}}
\renewcommand{\euro}{\ifmmode\text{\texteuro}\else\texteuro\fi}
\DeclareUnicodeCharacter{00B0}{\ifmmode^{\circ}\else\textdegree\fi}
\DeclareUnicodeCharacter{20AC}{\ifmmode\text{\texteuro}\else\texteuro\fi}
% Énoncés : Éducation nationale. Transcription et corrigés : APMEP, Denis Vergès et contributeurs. Découpage et classement : CoopMaths, Sébastien Lozano, Éric Elter et contributeurs. Adaptation de la présentation : Maths Explorer.
% CC BY-SA : https://coopmaths.fr/www/about/
% https://coopmaths.fr/alea/static/bac/2025/tex/bac_2025_06_sujet1_etrangers_1_cor.tex
\begin{document}
\noindent\textit{APMEP / CoopMaths. CC BY-SA. Présentation adaptée par Maths Explorer.}\par\medskip
\medskip

\textbf{Partie A : }

\
\begin{enumerate}
\item Un arbre pondéré représentant la situation est : 

\hfill~
\begin{tikzpicture}[xscale=1,yscale=1,baseline={(R.base)}]
\tikzstyle{fleche}=[->,>=latex,thick]
\tikzstyle{noeud}=[fill=white,circle,inner sep=2pt]
\tikzstyle{feuille}=[fill=white,circle,inner sep=2pt]
\tikzstyle{etiquette}=[pos=0.6,fill=white, inner xsep=3pt, inner ysep=1.5pt]
\def\DistanceInterNiveaux{3}
\def\DistanceInterFeuilles{0.9}
\def\NiveauA{(0)*\DistanceInterNiveaux}
\def\NiveauB{(1)*\DistanceInterNiveaux}
\def\NiveauC{(2)*\DistanceInterNiveaux}
\def\InterFeuilles{(-1)*\DistanceInterFeuilles}

\node[noeud] (R) at ({\NiveauA},{(1.5)*\InterFeuilles}) {};
\node[noeud] (Ra) at ({\NiveauB},{(0.5)*\InterFeuilles}) {$T$};
\node[feuille] (Raa) at ({\NiveauC},{(0)*\InterFeuilles}) {$E$};
\node[feuille] (Rab) at ({\NiveauC},{(1)*\InterFeuilles}) {$\overline{E}$};
\node[noeud] (Rb) at ({\NiveauB},{(2.5)*\InterFeuilles}) {$\overline{T}$};
\node[feuille] (Rba) at ({\NiveauC},{(2)*\InterFeuilles}) {$E$};
\node[feuille] (Rbb) at ({\NiveauC},{(3)*\InterFeuilles}) {$\overline{E}$};

\draw[fleche] (R.east)--(Ra.west) node[etiquette] {$0,1$};
\draw[fleche] (Ra.east)--(Raa.west) node[etiquette] {$0,3$};
\draw[fleche] (Ra.east)--(Rab.west) node[etiquette] {$0,7$};
\draw[fleche] (R.east)--(Rb.west) node[etiquette] {$0,9$};
\draw[fleche] (Rb.east)--(Rba.west) node[etiquette] {$0,15$};
\draw[fleche] (Rb.east)--(Rbb.west) node[etiquette] {$0,85$};
\end{tikzpicture} \hfill~

$P\left(\overline{T} \cap E\right) = P\left(\overline{T} \right)\times P_{\overline{T}}(E)= 0,9 \times 0,15 = 0,135$.
\item Les évènements $T$ et $\overline{T}$ forment une partition de l'univers, donc, d'après la formule des probabilités totales : \quad $P(E) = P\big(\,T\cap E\big) + P\big(\,\overline{T}\cap E\big) $

On a donc : \quad $P(E) = 0,1 \times 0,3 + 0,135 = 0,165$.

La probabilité qu'une erreur soit détectée lors du contrôle est égale à $0,165$.

\item $P_{E}(T)= \dfrac{P(T\cap E)}{P(E)} = \dfrac{0,1 \times 0,3}{0,165}\approx \np{0,1818}$.

La probabilité qu'un contrôle total ait été effectué, sachant qu'une erreur a été détectée vaut  environ $0,18$, arrondie au centième.
\end{enumerate}

\textbf{Partie B :}
Sur une journée donnée, une caisse automatique déclenche 15 contrôles. La probabilité qu'un contrôle mette en évidence une erreur est $p=0,165$. La détection d'une erreur lors d'un contrôle est indépendante des autres contrôles.

On note $X$ la variable aléatoire égale au nombre d'erreurs détectées lors des contrôles de cette journée.

\begin{enumerate}
\item Sur une journée donnée, une caisse automatique déclenche 15 contrôles. La probabilité qu'un contrôle mette en évidence une erreur est $p=0,165$ donc $X$ suit une loi binomiale de paramètres $n=15$ et $p=0,165$. 
\item $\begin{aligned}[t] P\left(X=5\right)
& = \displaystyle \binom{15}{5}\times 0,165^{5}\times 0,835^{15 - 5}\\
& = \np{3003}\times 0,165^{5} \times 0,835^{10} \\
& \approx \np{0,0605}
\end{aligned}$.

Finalement, avec les consignes d'arrondi : \quad $P\left(X=5\right) \approx 0,06$

La probabilité qu'exactement 5 erreurs soient détectées vaut $0,06$.

\item  $P(X\geqslant 1)=1-P(X=0)= 1- 0,835^
{15} \approx \np{0,9331}$.

La probabilité qu'au moins une erreur soit détectée vaut $0,93$.

\item Soit $n$ le nombre de contrôles effectués chaque jour.

Soit $Y$ la variable aléatoire égale au nombre d'erreurs détectées lors des contrôles de cette journée.

$Y$ suit une loi binomiale de paramètres $n$ et $p=0,165$.

$P(Y\geqslant 1)=1-P(Y=0)= 1- 0,835^n$

On souhaite que la probabilité qu'au moins une erreur soit détectée chaque jour soit supérieure à 99 \%.

On veut donc que $1- 0,835^n \geqslant 0,99$.

$\begin{aligned} 1- 0,835^n \geqslant 0,99 & \iff 0,01 \geqslant 0,835^n\\
& \iff \ln(0,01)\geqslant \ln \left(0,835^n\right) \, \text{car la fonction logarithme est strictement}\\
& \qquad \text{croissante sur } ]0~;~+\infty[\\
 &\iff \ln (0,01)\geqslant n \ln \left(0,835\right)\\
& \iff \dfrac{\ln(0,01)}{\ln(0,835)} \leqslant n \text{ car } \ln(0,935)<0\\
\end{aligned}$

or $ \dfrac{\ln(0,01)}{\ln(0,835)} \approx 25,5$

Il faut déclencher 26 contrôles chaque jour.
\end{enumerate}

\textbf{Partie C : }

Le magasin comporte trois caisses automatiques identiques qui, lors d'une journée, ont chacune déclenché 20 contrôles. On note $X_{1}, X_{2}$ et $X_{3}$ les variables aléatoires associant à chacune des caisses le nombre d'erreurs détectées lors de cette journée.

On admet que les variables aléatoires $X_{1}, X_{2}$ et $X_{3}$ sont indépendantes entre elles et suivent chacune une loi binomiale $\mathcal{B}(20~;~0,165)$.

\begin{enumerate}
\item  $X_1$ suit une loi binomiale de paramètres $n=20$ et $p=0,175$. 

Donc $E\left( X_1\right) = n\times p = 20 \times 0,165 = 3,3$ 

et $V\left( X_1\right) = n\times p \times (1-p) = 20 \times 0,165 \times 0,835 = \np{2,7555}$ 

\item $E(S)=E\left(X_{1}+X_{2}+X_{3}\right) = E\left( X_1\right) +E\left( X_2\right)+E\left( X_3\right) = 3 \times E\left( X_1\right) = 3 \times 3,3 = 9,9$ 

Les variables $X_{1}, X_{2}$ et $X_{3}$ sont indépendantes donc : 

$V(S)=V\left(X_{1}+X_{2}+X_{3}\right) = V\left( X_1\right) +V\left( X_2\right)+V\left( X_3\right) = 3 \times V\left( X_1\right) = 3 \times 2,7555 =\np{8,2665}$.
\item L'évènement \og $6 <S< 14$ \fg{} revient à $\big|S - 10\big| < 4$ et

$P\Big(\big|S- 10 \big| < 4\Big) = 1 - P\Big(\big|S - 10 \big| \geqslant 4\Big)$.

Or, pour tout réel $t>0$, d'après l'inégalité de Bienaymé-Tchebychev, on a :

$P\Big(\big|S-E\big(\,S\big)\big| \geqslant t\Big) \leqslant \dfrac{V\big(\,S\big)}{t^2}$

Pour $t = 4$ on obtient : $P\Big(\big|S-10\big| \geqslant 4\Big)\leqslant \dfrac{8,2665}{4^2}$.

D'où $P\Big(\big|S- 10 \big| < 4\Big) = 1 - P\Big(\big|S - 10 \big| \geqslant 4\Big) \geqslant 1- \dfrac{8,2665}{16}$.

Or $1 - \dfrac{8,2665}{16} \approx \np{0,4833}$

donc la probabilité que le nombre total d'erreurs sur la journée soit strictement compris entre 6 et 14 est bien supérieure à $0,48$.
\end{enumerate}
\medskip

\end{document}
